\scanpage{86}
\mspage{7}

groupoïdes envisagé est représentable par le schéma

\[
U_{03}\times \mathord{\mathrm{I}\mkern-7mu\mu}_4^*
\tag*{(37)}
\]

Le fibré $\Delta_E$ correspondant à la courbe elliptique épinglée (par Legendre-Torelli) universelle sur la base $U_{03}\times \mathord{\mathrm{I}\mkern-7mu\mu}_4^*$, est isomorphe canoniquement à

\[
\begin{aligned}
&\deleted{\underline O_{U_{0,3}}(1)^{\unsure{'}}}\\
&\deleted{T_0=\deleted{\unsure{U_{0,3}^{\natural}|U_{0,3}}}\underline O_{U_{0,3}}(1)^\natural\times \mathord{\mathrm{I}\mkern-7mu\mu}}
\end{aligned}
\]

\[
\mathbb{T}_0=\underline O_{U_{0,3}\times \mathord{\mathrm{I}\mkern-7mu\mu}_4^*}(1)^{\natural}
\tag*{(38)}
\]

où ici, puisqu'on est \deleted{sur} \inserted{au dessus} $\mu_4^*$, le fibré en carrés $Q_{0S}$ $(S=U_{03}\times \mathord{\mathrm{I}\mkern-7mu\mu}_4^*)$ est trivial, et on a un iso. can.

\[
\mathbb{T}_0\simeq\underline O_{U_{0,3}}(1)
\tag*{(39)}
\]

\deleted{\unsure{Le groupe|La question} modulaire pour la situation, est le groupe} \inserted{Mais il faut expliciter le ``groupe modulaire'' naturel pour cette situation de} structures épinglées, et examiner son action sur $U_{03}\times \mathord{\mathrm{I}\mkern-7mu\mu}_4^*$ ainsi que sur $\mathbb{T}_0$.

% archive sheet number: 85
